%%%%%%%%%%%%%%%%%%%%%%%%%%%%%%%%%%%%%%%%%%%%%%%%%%%%%%%%%%%%%%%%%%%%%%%%%%%%%%%
% Harvard Academic Notes - 통합 마스터 템플릿
% 모든 강의 노트에 적용되는 통일된 스타일
% 버전: 2.1 - 가독성 개선 (선택적 최적화)
% 최종 수정일: 2025-11-17
%%%%%%%%%%%%%%%%%%%%%%%%%%%%%%%%%%%%%%%%%%%%%%%%%%%%%%%%%%%%%%%%%%%%%%%%%%%%%%%

\documentclass[11pt,a4paper]{article}

%========================================================================================
% 기본 패키지
%========================================================================================

% --- 한국어 지원 ---
\usepackage{kotex}

% --- 페이지 레이아웃 ---
\usepackage[top=20mm, bottom=20mm, left=20mm, right=18mm]{geometry}
\usepackage{setspace}
\onehalfspacing                      % 1.5배 줄간격
\setlength{\parskip}{0.5em}          % 문단 간격
\setlength{\parindent}{0pt}          % 들여쓰기 없음

% --- 표 관련 ---
\usepackage{booktabs}              % 고품질 표
\usepackage{tabularx}              % 자동 너비 조절 표
\usepackage{array}                 % 표 컬럼 확장
\usepackage{longtable}             % 여러 페이지 표
\renewcommand{\arraystretch}{1.1}  % 표 행간 조절

%========================================================================================
% 헤더 및 푸터
%========================================================================================

\usepackage{fancyhdr}
\pagestyle{fancy}
\fancyhf{}
\fancyhead[L]{\small\textit{CSCI E-103: Introduction to the Intellectual Enterprises of CS}}
\fancyhead[R]{\small\textit{Module 06}}
\fancyfoot[C]{\thepage}
\renewcommand{\headrulewidth}{0.5pt}
\renewcommand{\footrulewidth}{0.3pt}

% 첫 페이지는 헤더 없음
\fancypagestyle{firstpage}{
    \fancyhf{}
    \fancyfoot[C]{\thepage}
    \renewcommand{\headrulewidth}{0pt}
}

%========================================================================================
% 색상 정의 (파스텔 톤 + 다크모드 호환)
%========================================================================================

\usepackage[dvipsnames]{xcolor}

% 밝은 배경용 파스텔 색상
\definecolor{lightblue}{RGB}{220, 235, 255}      % 부드러운 파랑
\definecolor{lightgreen}{RGB}{220, 255, 235}     % 부드러운 초록
\definecolor{lightyellow}{RGB}{255, 250, 220}    % 부드러운 노랑
\definecolor{lightpurple}{RGB}{240, 230, 255}    % 부드러운 보라
\definecolor{lightgray}{gray}{0.95}              % 밝은 회색
\definecolor{lightpink}{RGB}{255, 235, 245}      % 부드러운 핑크
\definecolor{boxgray}{gray}{0.95}
\definecolor{boxblue}{rgb}{0.9, 0.95, 1.0}
\definecolor{boxred}{rgb}{1.0, 0.95, 0.95}
\definecolor{myblue}{RGB}{50, 100, 180}

% 진한 색상 (테두리/제목용)
\definecolor{darkblue}{RGB}{50, 80, 150}
\definecolor{darkgreen}{RGB}{40, 120, 70}
\definecolor{darkorange}{RGB}{200, 100, 30}
\definecolor{darkpurple}{RGB}{100, 60, 150}

%========================================================================================
% 박스 환경 (tcolorbox) - 6가지 타입
%========================================================================================

\usepackage[most]{tcolorbox}
\tcbuselibrary{skins, breakable}

% 1. 개요 박스 (강의 시작 부분)
\newtcolorbox{overviewbox}[1][]{
    enhanced,
    colback=lightpurple,
    colframe=darkpurple,
    fonttitle=\bfseries\large,
    title=📚 강의 개요,
    arc=3mm,
    boxrule=1pt,
    left=8pt,
    right=8pt,
    top=8pt,
    bottom=8pt,
    breakable,
    #1
}

% 2. 요약 박스
\newtcolorbox{summarybox}[1][]{
    enhanced,
    colback=lightblue,
    colframe=darkblue,
    fonttitle=\bfseries,
    title=📝 핵심 요약,
    arc=2mm,
    boxrule=0.7pt,
    left=6pt,
    right=6pt,
    top=6pt,
    bottom=6pt,
    breakable,
    #1
}

% 3. 핵심 정보 박스
\newtcolorbox{infobox}[1][]{
    enhanced,
    colback=lightgreen,
    colframe=darkgreen,
    fonttitle=\bfseries,
    title=💡 핵심 정보,
    arc=2mm,
    boxrule=0.7pt,
    left=6pt,
    right=6pt,
    top=6pt,
    bottom=6pt,
    breakable,
    #1
}

% 4. 주의사항 박스
\newtcolorbox{warningbox}[1][]{
    enhanced,
    colback=lightyellow,
    colframe=darkorange,
    fonttitle=\bfseries,
    title=⚠️ 주의사항,
    arc=2mm,
    boxrule=0.7pt,
    left=6pt,
    right=6pt,
    top=6pt,
    bottom=6pt,
    breakable,
    #1
}

% 5. 예제 박스
\newtcolorbox{examplebox}[1][]{
    enhanced,
    colback=lightgray,
    colframe=black!60,
    fonttitle=\bfseries,
    title=📖 예제,
    arc=2mm,
    boxrule=0.7pt,
    left=6pt,
    right=6pt,
    top=6pt,
    bottom=6pt,
    breakable,
    #1
}

% 6. 정의 박스
\newtcolorbox{definitionbox}[1][]{
    enhanced,
    colback=lightpink,
    colframe=purple!70!black,
    fonttitle=\bfseries,
    title=📌 정의,
    arc=2mm,
    boxrule=0.7pt,
    left=6pt,
    right=6pt,
    top=6pt,
    bottom=6pt,
    breakable,
    #1
}

% 7. 중요 박스 (importantbox - warningbox와 유사)
\newtcolorbox{importantbox}[1][]{
    enhanced,
    colback=boxred,
    colframe=red!70!black,
    fonttitle=\bfseries,
    title=⚠️ 매우 중요,
    arc=2mm,
    boxrule=0.7pt,
    left=6pt,
    right=6pt,
    top=6pt,
    bottom=6pt,
    breakable,
    #1
}

% 8. cautionbox (warningbox와 동일)
\let\cautionbox\warningbox
\let\endcautionbox\endwarningbox

% tcbset 스타일 정의 (\begin{tcolorbox}[summarybox] 형식 사용 시 호환)
\tcbset{
    summarybox/.style={enhanced, colback=lightblue, colframe=darkblue, fonttitle=\bfseries, title=핵심 요약, arc=2mm, boxrule=0.7pt, breakable},
    warningbox/.style={enhanced, colback=lightyellow, colframe=darkorange, fonttitle=\bfseries, title=주의, arc=2mm, boxrule=0.7pt, breakable},
    examplebox/.style={enhanced, colback=lightgray, colframe=black!60, fonttitle=\bfseries, title=예제, arc=2mm, boxrule=0.7pt, breakable},
    definitionbox/.style={enhanced, colback=lightpink, colframe=purple!70!black, fonttitle=\bfseries, title=정의, arc=2mm, boxrule=0.7pt, breakable},
    importantbox/.style={enhanced, colback=boxred, colframe=red!70!black, fonttitle=\bfseries, title=매우 중요, arc=2mm, boxrule=0.7pt, breakable},
    infobox/.style={enhanced, colback=lightgreen, colframe=darkgreen, fonttitle=\bfseries, title=핵심 정보, arc=2mm, boxrule=0.7pt, breakable},
    conceptbox/.style={enhanced, colback=lightgreen, colframe=darkgreen, fonttitle=\bfseries, title=개념, arc=2mm, boxrule=0.7pt, breakable},
    analogybox/.style={enhanced, colback=green!5!white, colframe=green!60!black, fonttitle=\bfseries, title=비유, arc=2mm, boxrule=0.7pt, breakable},
    overviewbox/.style={enhanced, colback=lightpurple, colframe=darkpurple, fonttitle=\bfseries\large, title=강의 개요, arc=3mm, boxrule=1pt, breakable},
    cautionbox/.style={enhanced, colback=lightyellow, colframe=darkorange, fonttitle=\bfseries, title=주의, arc=2mm, boxrule=0.7pt, breakable},
    mybox/.style={enhanced, colback=lightblue, colframe=darkblue, fonttitle=\bfseries, title={#1}, arc=2mm, boxrule=0.7pt, breakable},
    definition/.style={enhanced, colback=lightpink, colframe=purple!70!black, fonttitle=\bfseries, title={#1}, arc=2mm, boxrule=0.7pt, breakable},
    example/.style={enhanced, colback=lightgray, colframe=black!60, fonttitle=\bfseries, title={#1}, arc=2mm, boxrule=0.7pt, breakable},
    summary/.style={enhanced, colback=lightblue, colframe=darkblue, fonttitle=\bfseries, title={#1}, arc=2mm, boxrule=0.7pt, breakable},
    warning/.style={enhanced, colback=lightyellow, colframe=darkorange, fonttitle=\bfseries, title={#1}, arc=2mm, boxrule=0.7pt, breakable},
}

%========================================================================================
% 코드 블록 설정 (밝은 배경)
%========================================================================================

\usepackage{listings}

\definecolor{codegray}{rgb}{0.5,0.5,0.5}
\definecolor{codepurple}{rgb}{0.58,0,0.82}
\definecolor{backcolour}{rgb}{0.95,0.95,0.95}

\lstset{
    basicstyle=\ttfamily\small,
    backgroundcolor=\color{lightgray},
    keywordstyle=\color{darkblue}\bfseries,
    commentstyle=\color{darkgreen}\itshape,
    stringstyle=\color{purple!80!black},
    numberstyle=\tiny\color{black!60},
    numbers=left,
    numbersep=8pt,
    breaklines=true,
    breakatwhitespace=false,
    frame=single,
    frameround=tttt,
    rulecolor=\color{black!30},
    captionpos=b,
    showstringspaces=false,
    tabsize=2,
    xleftmargin=15pt,
    xrightmargin=5pt,
    escapeinside={\%*}{*)}
}

% Python 코드 스타일
\lstdefinestyle{pythonstyle}{
    language=Python,
    morekeywords={self, True, False, None},
}

% SQL 코드 스타일
\lstdefinestyle{sqlstyle}{
    language=SQL,
    morekeywords={SELECT, FROM, WHERE, JOIN, GROUP, BY, ORDER, HAVING},
}

%========================================================================================
% 목차 스타일링
%========================================================================================

\usepackage{tocloft}
\renewcommand{\cftsecleader}{\cftdotfill{\cftdotsep}}
\setlength{\cftbeforesecskip}{0.4em}
\renewcommand{\cftsecfont}{\bfseries}
\renewcommand{\cftsubsecfont}{\normalfont}

%========================================================================================
% 표 및 그림
%========================================================================================

\usepackage{graphicx}              % 이미지
\usepackage{adjustbox}             % 표/박스 크기 조절

% 표 캡션 스타일
\usepackage{caption}
\captionsetup[table]{
    labelfont=bf,
    textfont=it,
    skip=5pt
}
\captionsetup[figure]{
    labelfont=bf,
    textfont=it,
    skip=5pt
}

%========================================================================================
% 수학
%========================================================================================

\usepackage{amsmath, amssymb, amsthm}

% 정리 환경
\theoremstyle{definition}
\newtheorem{theorem}{정리}[section]
\newtheorem{lemma}[theorem]{보조정리}
\newtheorem{proposition}[theorem]{명제}
\newtheorem{corollary}[theorem]{따름정리}
\newtheorem{definition}{정의}[section]
\newtheorem{example}{예제}[section]

%========================================================================================
% 하이퍼링크
%========================================================================================

\usepackage[
    colorlinks=true,
    linkcolor=blue!80!black,
    urlcolor=blue!80!black,
    citecolor=green!60!black,
    bookmarks=true,
    bookmarksnumbered=true,
    pdfborder={0 0 0}
]{hyperref}

% PDF 메타데이터는 각 문서에서 설정
\hypersetup{
    pdftitle={CSCI E-103: Introduction to the Intellectual Enterprises of CS - Module 06},
    pdfauthor={강의 노트},
    pdfsubject={Academic Notes}
}

%========================================================================================
% 기타 유용한 패키지
%========================================================================================

\usepackage{enumitem}              % 리스트 커스터마이징
\setlist{nosep, leftmargin=*, itemsep=0.3em}

\usepackage{microtype}             % 타이포그래피 개선
\usepackage{footnote}              % 각주 개선
\usepackage{url}                   % URL 줄바꿈
\urlstyle{same}

%========================================================================================
% 사용자 정의 명령어
%========================================================================================

% 강조 텍스트
\newcommand{\important}[1]{\textbf{\textcolor{red!70!black}{#1}}}
\newcommand{\keyword}[1]{\textbf{#1}}
\newcommand{\term}[1]{\textit{#1}}
\newcommand{\code}[1]{\texttt{#1}}

% 용어 설명 (인라인)
\newcommand{\defterm}[2]{\textbf{#1}\footnote{#2}}

% 섹션 시작 전 페이지 분리
\newcommand{\newsection}[1]{\newpage\section{#1}}

%========================================================================================
% 문서 제목 스타일
%========================================================================================

\usepackage{titling}
\pretitle{\begin{center}\LARGE\bfseries}
\posttitle{\par\end{center}\vskip 0.5em}
\preauthor{\begin{center}\large}
\postauthor{\end{center}}
\predate{\begin{center}\large}
\postdate{\par\end{center}}

%========================================================================================
% 섹션 제목 간격
%========================================================================================

\usepackage{titlesec}
\titlespacing*{\section}{0pt}{1.5em}{0.8em}
\titlespacing*{\subsection}{0pt}{1.2em}{0.6em}
\titlespacing*{\subsubsection}{0pt}{1em}{0.5em}

%========================================================================================
% 메타 정보 박스 명령어
%========================================================================================

\newcommand{\metainfo}[4]{
\begin{tcolorbox}[
    colback=lightpurple,
    colframe=darkpurple,
    boxrule=1pt,
    arc=2mm,
    left=10pt,
    right=10pt,
    top=8pt,
    bottom=8pt
]
\begin{tabular}{@{}rl@{}}
▣ \textbf{강의명:} & #1 \\[0.3em]
▣ \textbf{주차:} & #2 \\[0.3em]
▣ \textbf{교수명:} & #3 \\[0.3em]
▣ \textbf{목적:} & \begin{minipage}[t]{0.75\textwidth}#4\end{minipage}
\end{tabular}
\end{tcolorbox}
}

%========================================================================================
% 끝
%========================================================================================


\begin{document}

\thispagestyle{firstpage}

\metainfo{CSCI103}{Lecture 6}{UIUC Student}{이 문서는 CSCI103 Lecture 6 강의 내용을 통합하여, 비즈니스 인텔리전스(BI) 분석 및 데이터 시각화의 핵심 개념을 다룹니다. 데이터 웨어하우스, 데이터 레이크, 그리고 이 둘의 장점을 결합한 레이크하우스 아키텍처의 진화를 이해하고, 메달리온 아키텍처, 데이터 메시, 데이터 패브릭과 같은 현대 데이터 관리 전략을 탐구합니다. 또한, Databricks 플랫폼의 BI 지원 기능, 주요 KPI(핵심 성과 지표), 데이터 모델링 기법, 그리고 BI 분석가에게 필요한 역량에 대해 상세히 설명합니다. 이 문서는 시험 대비를 위한 완벽한 학습 자료로, 모든 개념을 초심자도 이해할 수 있도록 친절하게 해설하고 풍부한 예시를 제공합니다.}

\tableofcontents

\newpage

\section{용어 정리}

다음 표는 강의에서 다루는 주요 용어들을 쉽게 설명하고 원어를 병기하여 이해를 돕습니다.

\begin{adjustbox}{width=\textwidth,center}
\begin{tabular}{lllc}
\toprule
\textbf{용어} & \textbf{쉬운 설명} & \textbf{원어} & \textbf{비고} \\
\midrule
BI & 기업이 데이터를 분석하여 더 나은 의사결정을 내리도록 돕는 과정과 기술 & Business Intelligence & '무엇이 일어났고, 어떻게 일어났는지' 설명 \\
BA & 과거 데이터를 바탕으로 미래를 예측하고 '왜' 일어났는지 설명하는 분석 & Business Analytics & '왜 일어났고, 다음에 무엇이 일어날지' 예측 \\
데이터 웨어하우스 & 정형 데이터를 저장하고 분석하기 위해 최적화된 시스템 & Data Warehouse & 구조화된 데이터, 빠른 검색, 배치 처리 \\
데이터 레이크 & 정형, 비정형 데이터를 모두 저장할 수 있는 대규모 저장소 & Data Lake & 모든 종류의 데이터, 저렴한 저장 비용, 스키마 온 리드 \\
레이크하우스 & 데이터 레이크의 유연성과 데이터 웨어하우스의 성능을 결합한 아키텍처 & Lakehouse & 데이터 레이크 위에 웨어하우스 기능 통합 \\
메달리온 아키텍처 & 데이터 품질과 정제 수준에 따라 데이터를 계층화하는 방식 (Bronze, Silver, Gold) & Medallion Architecture & 데이터 정제 및 분석 준비 단계 \\
데이터 사일로 & 데이터가 조직 내 여러 팀이나 시스템에 분리되어 연결되지 않은 상태 & Data Silo & 데이터 통합의 어려움 초래 \\
데이터 스웜프 & 데이터 레이크에 데이터가 무분별하게 쌓여 관리가 안 되는 상태 & Data Swamp & 데이터 레이크의 실패 사례 \\
KPI & 비즈니스 목표 달성 정도를 측정하는 핵심 지표 & Key Performance Indicator & 비즈니스 성과 모니터링 \\
ETL & 데이터 추출, 변환, 로드 과정 & Extract, Transform, Load & 데이터 웨어하우스 구축의 핵심 과정 \\
페더레이션 & 데이터를 물리적으로 이동시키지 않고 여러 소스에서 통합하여 쿼리하는 방식 & Federation & 데이터 소유권은 원본에 유지, 읽기 전용 \\
Unity Catalog & Databricks의 통합 거버넌스 솔루션 & Unity Catalog & 데이터, 모델, 피처에 대한 접근 제어, 계보 관리 \\
Databricks SQL & Databricks 플랫폼에서 BI 및 데이터 웨어하우징을 위한 SQL 인터페이스 & Databricks SQL (DBSQL) & 고성능 SQL 쿼리 및 대시보드 지원 \\
Photon & Databricks의 고성능 쿼리 엔진 (C++로 재작성된 Spark) & Photon & Spark 워크로드 가속화 \\
Genie & Databricks 대시보드에서 자연어로 질문하고 답변을 얻는 AI 기능 & Genie & 비즈니스 사용자의 셀프 서비스 BI 지원 \\
AI 쿼리/함수 & SQL 쿼리 내에서 LLM 모델을 호출하여 데이터 분석을 강화하는 기능 & AI Query/Function & 감성 분석, 분류, 정보 추출 등 \\
\bottomrule
\end{tabular}
\end{adjustbox}

\newpage
\section{핵심 개념·원리}

\subsection{데이터 저장소의 진화}
데이터 저장 및 분석 기술은 지난 수십 년간 비즈니스 요구사항에 맞춰 끊임없이 발전해 왔습니다. 이 진화는 데이터의 양, 다양성, 속도 증가에 대응하고, 더 효율적이고 통찰력 있는 분석을 가능하게 하는 방향으로 이루어졌습니다.

\begin{overviewbox}[title=데이터 저장소 진화의 여정]
데이터 저장소의 진화는 단순히 기술의 변화를 넘어, 기업이 데이터를 통해 가치를 창출하는 방식의 변화를 의미합니다. 초기에는 수동적인 데이터 관리에서 시작하여, 이제는 AI와 자연어 처리를 통해 누구나 데이터에 접근하고 통찰력을 얻을 수 있는 시대로 나아가고 있습니다.
\end{overviewbox}

\begin{itemize}
    \item \textbf{스프레드시트 (초기):} 데이터 수집 및 분석의 초기 형태로, 소규모 데이터 관리에 사용되었습니다.
    \item \textbf{데이터 웨어하우스 (Data Warehouse, DW):}
    \begin{itemize}
        \item Bill Inmon과 Ralph Kimball에 의해 발전.
        \item Inmon은 ER(개체-관계) 모델에서 차원 모델을 생성하는 방식을, Kimball은 ER 단계를 건너뛰고 바로 팩트(Fact) 및 차원(Dimensional) 테이블(스타 또는 스노우플레이크 스키마)을 사용하는 방식을 제안했습니다.
    \end{itemize}
    \item \textbf{대규모 병렬 처리(MPP) 데이터베이스 (Greenplum, Teradata):}
    \begin{itemize}
        \item 데이터 양이 너무 커져 기존 DW에 맞지 않게 되면서 등장.
        \item 독점 시스템을 사용하여 데이터와 컴퓨팅을 여러 머신에 분산시켜 테라바이트(Terabyte) 단위의 데이터를 처리할 수 있게 했습니다.
        \item 하지만 매우 비싸고 비용 효율적이지 못했습니다.
    \end{itemize}
    \item \textbf{NoSQL 솔루션 (Google Bigtable):}
    \begin{itemize}
        \item Google의 Bigtable은 테라바이트에서 페타바이트(Petabyte)까지 처리할 수 있는 컬럼형 데이터베이스로, NoSQL(Not Only SQL) 솔루션의 시작을 알렸습니다.
        \item 이는 대규모 테이블형 정보에 대한 쿼리를 가능하게 했습니다.
    \end{itemize}
    \item \textbf{Hadoop (Mike Carett Caparella, Doug Cutting):}
    \begin{itemize}
        \item 저렴한 하드웨어 플랫폼에서 엄청난 양의 데이터를 저장하고 처리할 수 있게 했습니다.
        \item 스토리지와 컴퓨팅을 분리하고 수평적으로 확장 가능한 분산 아키텍처를 사용했습니다.
    \end{itemize}
    \item \textbf{데이터 레이크 (Data Lake):}
    \begin{itemize}
        \item Hadoop의 등장으로 데이터 레이크가 탄생했습니다.
        \item 그러나 적절한 거버넌스 없이는 쉽게 데이터 스웜프(Data Swamp)로 변질될 수 있었습니다.
    \end{itemize}
    \item \textbf{데이터 메시 및 패브릭 (Data Mesh & Fabric):}
    \begin{itemize}
        \item 데이터가 서비스 또는 제품이 되는 개념으로, 데이터 거버넌스와 접근성을 향상시켰습니다.
    \end{itemize}
    \item \textbf{레이크하우스 (Lakehouse):}
    \begin{itemize}
        \item 데이터 레이크의 유연성과 데이터 웨어하우스의 기능을 결합하여, 모든 데이터를 함께 유지하고 좋은 거버넌스를 지원합니다.
        \item 기본 데이터 레이크 플랫폼에서 직접 데이터 웨어하우스 기능을 지원합니다.
    \end{itemize}
\end{itemize}

\begin{warningbox}[title=지속적인 진화]
이러한 데이터 기술의 진화는 아직 끝나지 않았습니다. AI, 특히 LLM(대규모 언어 모델)의 발전은 미래에 더 많은 변화를 가속화할 것으로 예상됩니다.
\end{warningbox}

\subsection{데이터 레이크, 데이터 웨어하우스, 레이크하우스 비교}
데이터를 저장하고 관리하는 세 가지 주요 아키텍처는 각각 고유한 장단점을 가지고 있습니다. 레이크하우스는 이들의 장점을 결합하여 현대 데이터 환경의 복잡성을 해결하고자 합니다.

\begin{summarybox}[title=핵심 요약: 레이크하우스의 등장 배경]
데이터 레이크는 유연하지만 성능과 거버넌스에 약점이 있었고, 데이터 웨어하우스는 성능이 좋지만 유연성과 확장성에 한계가 있었습니다. 레이크하우스는 이 두 가지 아키텍처의 장점을 결합하여, 대규모의 다양한 데이터를 고성능으로 처리하면서도 강력한 거버넌스를 제공하는 새로운 패러다임입니다.
\end{summarybox}

\begin{table}[h!]
\centering
\caption{데이터 레이크, 데이터 웨어하우스, 레이크하우스 비교}
\label{tab:data_storage_comparison}
\begin{adjustbox}{width=\textwidth,center}
\begin{tabular}{p{0.15\textwidth}p{0.25\textwidth}p{0.25\textwidth}p{0.25\textwidth}}
\toprule
\textbf{특징} & \textbf{데이터 레이크 (Data Lake)} & \textbf{데이터 웨어하우스 (Data Warehouse)} & \textbf{레이크하우스 (Lakehouse)} \\
\midrule
\textbf{데이터 유형} & 모든 유형 (정형, 반정형, 비정형) & 주로 정형 데이터 & 모든 유형 (정형, 반정형, 비정형) \\
\textbf{파일 형식} & 개방형 파일 형식 (Parquet, ORC, JSON 등) & 독점 형식 또는 특정 DB 형식 & 개방형 파일 형식 (Delta Lake) \\
\textbf{스키마} & 스키마 온 리드 (Schema on Read) & 스키마 온 라이트 (Schema on Write) & 스키마 온 리드/라이트 모두 지원 \\
\textbf{스토리지} & 저렴한 객체 스토리지 (S3, ADLS) & 고비용의 전용 스토리지 & 저렴한 객체 스토리지 \\
\textbf{컴퓨팅} & 스토리지와 분리, 효율적 사용 & 스토리지와 결합 또는 분리 & 스토리지와 분리, 효율적 사용 \\
\textbf{성능} & 대규모 데이터 처리, 유연성 높음 & 고성능 쿼리, 낮은 지연 시간 & 고성능 쿼리, 낮은 지연 시간 (DW 수준) \\
\textbf{접근 제어} & 세분화된 접근 제어 어려움 & 세분화된 접근 제어 용이 & 세분화된 접근 제어 용이 (DW 수준) \\
\textbf{사용자} & 데이터 과학자, ML 엔지니어 & 비즈니스 사용자, BI 분석가 & 모든 사용자 (BI, DS, ML) \\
\textbf{주요 용도} & 데이터 탐색, ML, 고급 분석 & BI 보고서, 대시보드 & BI, ML, 고급 분석, 실시간 처리 \\
\textbf{단점} & 거버넌스 부족 시 데이터 스웜프화 & 데이터 양 제한, 비정형 데이터 처리 어려움, 비용 & (이전 단점들을 해결) \\
\bottomrule
\end{tabular}
\end{adjustbox}
\end{table}

\subsubsection{데이터 웨어하우스 아키텍처}
\begin{itemize}
    \item \textbf{구조:} 정형 데이터가 ETL(Extract, Transform, Load) 프로세스를 통해 시스템에 공급됩니다.
    \item \textbf{스키마:} '스키마 온 라이트(Schema on Write)' 방식을 사용하여 데이터가 저장되기 전에 스키마가 정의됩니다.
    \item \textbf{구조:} 차원 모델과 같은 구조로 저장될 수 있으며, 보고서 및 BI 도구를 통해 노출됩니다.
\end{itemize}

\subsubsection{데이터 레이크 아키텍처}
\begin{itemize}
    \item \textbf{구조:} 데이터 레이크는 모든 종류의 데이터를 저장하며, 여기서 ETL 프로세스를 사용하여 데이터 웨어하우스를 채울 수 있습니다.
    \item \textbf{활용:} 보고서, BI, 데이터 과학, 머신러닝 등을 지원합니다.
    \item \textbf{장점:} 대규모 데이터와 정형/비정형 데이터를 모두 지원할 수 있어 더 유연합니다.
    \item \textbf{단점:} 데이터 웨어하우스가 여전히 별도로 존재해야 하므로 다소 번거로울 수 있습니다.
\end{itemize}

\subsubsection{레이크하우스 아키텍처}
\begin{itemize}
    \item \textbf{구조:} 데이터 웨어하우스 기능이 데이터 레이크의 일부로 포함됩니다.
    \item \textbf{기반:} 데이터 레이크가 구조화된, 반정형, 비정형 데이터를 모두 저장하는 기반 역할을 합니다.
    \item \textbf{활용:} 레이크하우스에서 직접 BI 보고서, 데이터 과학, 머신러닝을 지원합니다.
    \item \textbf{장점:} 별도의 데이터 웨어하우스에 투자할 필요가 없어 더 간단하고, 관리하기 쉬우며, 비용 효율적입니다.
    \item \textbf{성능:} 기본적으로 레이크 위에 구축되므로, 데이터 웨어하우스의 성능을 유지하면서도 훨씬 더 많은 양의 데이터를 처리할 수 있습니다.
\end{itemize}

\begin{examplebox}[title=레이크하우스의 장점]
레이크하우스는 마치 '만능 칼'과 같습니다. 데이터 레이크처럼 모든 종류의 재료(데이터)를 보관할 수 있고, 데이터 웨어하우스처럼 정교한 요리(분석)를 빠르게 만들 수 있는 도구(SQL 쿼리)도 내장하고 있습니다. 덕분에 여러 도구를 왔다 갔다 할 필요 없이 한 곳에서 모든 작업을 처리할 수 있어 시간과 비용을 절약할 수 있습니다.
\end{examplebox}

\subsection{메달리온 아키텍처 (Medallion Architecture)}
메달리온 아키텍처는 데이터 레이크 내에서 데이터의 품질과 정제 수준에 따라 데이터를 계층화하는 방식입니다. 이는 데이터 거버넌스를 강화하고, 데이터의 재사용성을 높이며, 오류 발생 시 복구 과정을 용이하게 합니다.

\begin{summarybox}[title=핵심 요약: 데이터 품질 관리]
메달리온 아키텍처는 데이터를 원본(Bronze), 정제(Silver), 분석 준비(Gold)의 세 가지 계층으로 나누어 관리함으로써, 데이터의 신뢰성을 높이고 다양한 사용자가 각자의 목적에 맞는 데이터를 효율적으로 사용할 수 있도록 돕습니다.
\end{summarybox}

\begin{itemize}
    \item \textbf{Bronze (브론즈) 계층:}
    \begin{itemize}
        \item \textbf{역할:} 원본(Raw) 데이터를 그대로 수집하여 저장하는 계층입니다.
        \item \textbf{특징:} 데이터가 수집된 형태 그대로 유지되며, 최소한의 변환만 적용됩니다. 오류가 있더라도 그대로 보존됩니다.
        \item \textbf{중요성:} 문제가 발생했을 때 언제든지 원본 데이터로 돌아가 다시 시작할 수 있는 '진실의 원천(Source of Truth)' 역할을 합니다.
    \end{itemize}
    \item \textbf{Silver (실버) 계층:}
    \begin{itemize}
        \item \textbf{역할:} 브론즈 계층의 데이터를 정제하고 통합하는 계층입니다.
        \item \textbf{특징:} 중복 제거, 데이터 클렌징, 조인, 데이터 형식 통일 등 기본적인 변환이 이루어집니다. 데이터 과학자나 분석가가 탐색적 분석을 수행하기에 적합한 형태입니다.
        \item \textbf{중요성:} 여러 소스의 데이터를 결합하여 일관된 뷰를 제공하고, 다음 단계의 분석을 위한 기반을 마련합니다.
    \end{itemize}
    \item \textbf{Gold (골드) 계층:}
    \begin{itemize}
        \item \textbf{역할:} 비즈니스 분석 및 보고서 작성을 위해 최적화된, 고도로 정제된 데이터를 저장하는 계층입니다.
        \item \textbf{특징:} 특정 비즈니스 요구사항에 맞춰 집계, 요약, 추가적인 변환이 적용됩니다. 비즈니스 사용자가 직접 BI 도구를 사용하여 통찰력을 얻을 수 있는 '분석 준비 완료(Analytics Ready)' 상태의 데이터입니다.
        \item \textbf{중요성:} 다양한 유형의 비즈니스 사용자(경영진, 마케터 등)와 활동(대시보드, 예측 모델)을 지원하며, 비즈니스 의사결정에 직접적으로 활용됩니다.
    \end{itemize}
\end{itemize}

\begin{examplebox}[title=올림픽 메달 비유]
메달리온 아키텍처는 올림픽 메달과 비슷합니다.
\begin{itemize}
    \item \textbf{브론즈(동메달):} 원본 데이터. 아직 다듬어지지 않은 원석 같은 상태입니다.
    \item \textbf{실버(은메달):} 정제된 데이터. 원석을 가공하여 어느 정도 형태를 갖춘 상태입니다. 데이터 과학자들이 실험하기 좋습니다.
    \item \textbf{골드(금메달):} 분석 준비 완료 데이터. 완벽하게 세공되어 바로 전시할 수 있는 보석과 같습니다. 비즈니스 의사결정에 바로 사용됩니다.
\end{itemize}
\end{examplebox}

\subsection{데이터 사일로 vs. 데이터 스웜프}
데이터 관리의 실패 사례를 보여주는 두 가지 개념입니다.

\begin{summarybox}[title=핵심 요약: 데이터 관리의 함정]
데이터 사일로는 데이터가 분리되어 소통이 안 되는 문제이고, 데이터 스웜프는 데이터가 너무 많고 무질서하여 쓸모없게 되는 문제입니다. 둘 다 데이터의 가치를 떨어뜨리고 비즈니스 의사결정을 방해합니다.
\end{summarybox}

\begin{itemize}
    \item \textbf{데이터 사일로 (Data Silo):}
    \begin{itemize}
        \item \textbf{정의:} 데이터가 조직 내의 여러 팀, 부서, 데이터베이스에 분리되어 저장되어 서로 연결되지 않는 상태를 의미합니다.
        \item \textbf{문제점:} 데이터가 고립되어 전체적인 비즈니스 뷰를 얻기 어렵고, 데이터 중복이 발생하며, 협업을 방해하고, 일관된 분석을 어렵게 만듭니다. 마치 각 부서가 자기만의 언어로 말하고 다른 부서의 말을 이해하지 못하는 것과 같습니다.
    \end{itemize}
    \item \textbf{데이터 스웜프 (Data Swamp):}
    \begin{itemize}
        \item \textbf{정의:} 데이터 레이크에 데이터가 무책임하게, 계획 없이 마구잡이로 덤프되어 관리되지 않는 상태를 의미합니다.
        \item \textbf{문제점:} 데이터에 대한 거버넌스(관리 규칙)가 없어 어떤 데이터가 어디에 있는지, 어떤 의미를 가지는지 알 수 없게 됩니다. 결국 데이터가 쓸모없게 되어 '데이터 늪'에 빠지는 것과 같습니다.
    \item \textbf{비유:} 데이터 레이크가 잘 관리된 도서관이라면, 데이터 스웜프는 책들이 아무렇게나 쌓여서 필요한 책을 찾을 수 없는 창고와 같습니다.
    \end{itemize}
\end{itemize}

\subsection{데이터 메시 vs. 데이터 패브릭}
현대 데이터 아키텍처에서 데이터의 접근성과 관리 효율성을 높이기 위한 두 가지 접근 방식입니다.

\begin{summarybox}[title=핵심 요약: 데이터 접근성 향상]
데이터 메시는 데이터를 제품처럼 취급하여 각 도메인 팀이 데이터를 소유하고 서비스하는 분산형 접근 방식이고, 데이터 패브릭은 기술 중심적으로 다양한 데이터 소스를 통합하여 데이터 접근을 단순화하는 방식입니다.
\end{summarybox}

\begin{table}[h!]
\centering
\caption{데이터 메시와 데이터 패브릭 비교}
\label{tab:data_mesh_fabric_comparison}
\begin{adjustbox}{width=\textwidth,center}
\begin{tabular}{p{0.15\textwidth}p{0.4\textwidth}p{0.4\textwidth}}
\toprule
\textbf{특징} & \textbf{데이터 메시 (Data Mesh)} & \textbf{데이터 패브릭 (Data Fabric)} \\
\midrule
\textbf{접근 방식} & 분산형, 도메인 중심 & 중앙 집중형, 기술 중심 \\
\textbf{데이터 소유권} & 각 도메인 팀이 데이터 제품을 소유 및 관리 & 데이터 패브릭 플랫폼이 데이터 통합 및 관리 \\
\textbf{주요 목표} & 데이터 민주화, 자율성, 확장성 & 데이터 통합, 접근 단순화, 자동화 \\
\textbf{데이터 처리} & 각 도메인에서 독립적으로 처리 & 패브릭 플랫폼을 통해 통합 처리 \\
\textbf{기술 초점} & 조직 문화, 거버넌스, 데이터 제품화 & 데이터 가상화, 그래프 기술, AI/ML 기반 자동화 \\
\textbf{예시} & 각 비즈니스 부서가 자체 데이터 API를 제공 & 여러 클라우드 및 온프레미스 데이터 소스를 단일 뷰로 통합 \\
\bottomrule
\end{tabular}
\end{adjustbox}
\end{table}

\begin{itemize}
    \item \textbf{데이터 메시 (Data Mesh):}
    \begin{itemize}
        \item \textbf{정의:} 데이터를 '제품(Product)'으로 취급하고, 분산된 도메인 중심 팀이 데이터를 소유하고 지원하는 아키텍처입니다.
        \item \textbf{핵심:} 각 팀은 자신들이 생성하거나 관리하는 데이터를 책임지고, 다른 팀이 쉽게 사용할 수 있도록 표준화된 인터페이스(API)를 통해 '데이터 제품'으로 노출합니다.
        \item \textbf{비유:} 마치 각 부서가 자신들의 전문 분야에 맞는 상품을 만들고, 이를 백화점(전사적 데이터 플랫폼)에 진열하여 다른 부서나 고객이 쉽게 구매할 수 있도록 하는 것과 같습니다.
    \end{itemize}
    \item \textbf{데이터 패브릭 (Data Fabric):}
    \begin{itemize}
        \item \textbf{정의:} 데이터가 어디에 있든(클라우드, 온프레미스, 엣지 등) 상관없이, 모든 데이터 소스를 통합하여 단일하고 일관된 데이터 뷰를 제공하는 기술 중심의 아키텍처입니다.
        \item \textbf{핵심:} 데이터 가상화, 그래프 기술, AI/ML 기반 자동화 등을 활용하여 데이터 통합, 거버넌스, 접근을 단순화합니다.
        \item \textbf{비유:} 마치 여러 언어를 사용하는 사람들이 모인 자리에서 실시간으로 모든 언어를 번역해주는 통역 시스템과 같습니다. 데이터가 어디에 있든 상관없이 필요한 정보를 바로 얻을 수 있도록 돕습니다.
    \end{itemize}
\end{itemize}

\subsection{비즈니스 인텔리전스(BI)의 이해}
비즈니스 인텔리전스(Business Intelligence, BI)는 기업이 데이터를 수집, 통합, 분석, 시각화하여 더 나은 비즈니스 의사결정을 내릴 수 있도록 돕는 일련의 기술과 애플리케이션, 프로세스를 의미합니다.

\begin{summarybox}[title=핵심 요약: 미래를 보는 수정구슬]
BI는 기업이 방대한 데이터를 통해 현재 상황을 이해하고, 미래를 예측하며, 경쟁 우위를 확보하기 위한 통찰력을 얻도록 돕는 '수정구슬'과 같습니다. 데이터 자체만으로는 부족하며, 데이터에서 추출된 '정보'가 비즈니스 성공의 열쇠입니다.
\end{summarybox}

\begin{itemize}
    \item \textbf{목표:} 비즈니스 의사결정을 개선하고, 비즈니스 성과를 향상시키는 것입니다.
    \item \textbf{활용 데이터:} 마케팅 데이터, 고객 데이터, 공급망 데이터, 자원 데이터 등 다양한 유형의 데이터를 활용합니다.
    \item \textbf{데이터와 정보:}
    \begin{itemize}
        \item \textbf{데이터 (Data):} 분석을 위해 필요한 원재료입니다.
        \item \textbf{정보 (Information):} 데이터에서 추출되어 비즈니스에 필요한 의미 있는 통찰력입니다. 데이터만으로는 충분하지 않으며, 정보가 있어야 비즈니스에서 성공할 수 있습니다.
    \end{itemize}
    \item \textbf{예시:}
    \begin{itemize}
        \item 고객 연락 및 상호작용 분석
        \item 영업 성사(Closed Deal) 분석
        \item Google Analytics를 통한 웹사이트 트래픽 분석
    \end{itemize}
\end{itemize}

\subsubsection{BI의 진화}
BI는 IT 부서의 통제에서 시작하여 점차 사용자 중심의 셀프 서비스 형태로 발전해 왔습니다.

\begin{itemize}
    \item \textbf{2000년대 초:} IT 부서가 데이터를 통제.
    \item \textbf{64비트 컴퓨터 도입:} 데이터 처리 능력 향상.
    \textbf{셀프 서비스 BI 및 클라우드 컴퓨팅 (AWS):} Amazon AWS의 등장으로 클라우드 기반 컴퓨팅이 확산되며 BI 접근성이 향상.
    \item \textbf{모바일 BI:} 모바일 기기를 통한 BI 접근 가능.
    \item \textbf{데이터 시각화 (PowerBI):} Microsoft PowerBI와 같은 도구로 데이터 시각화가 보편화.
    \item \textbf{머신러닝 및 AI:} BI에 머신러닝과 AI가 통합되어 예측 및 추천 기능 강화.
    \item \textbf{로우 코드/노 코드 (Low Code/No Code):} 비전문가도 BI 도구를 쉽게 사용할 수 있게 됨.
    \item \textbf{현재 (ChatGPT 등 LLM):} 자연어 대화형 인터페이스를 통해 누구나 BI에 접근할 수 있게 되어 'BI 민주화' 가속화.
\end{itemize}

\begin{warningbox}[title=데이터 엔지니어링의 중요성]
LLM을 통한 BI 민주화가 가속화되고 있지만, 이러한 도구들이 사용할 수 있도록 데이터를 준비하는 '데이터 엔지니어링'의 중요성은 여전히 매우 큽니다.
\end{warningbox}

\subsection{BI와 BA의 차이}
BI와 BA는 모두 데이터를 활용하여 비즈니스 가치를 창출하지만, 초점과 질문의 유형에서 차이가 있습니다.

\begin{summarybox}[title=핵심 요약: 과거 vs. 미래]
BI는 주로 '무엇이 일어났고, 어떻게 일어났는지'를 설명하는 과거 지향적 분석(기술적 분석)에 초점을 맞춥니다. 반면 BA는 '왜 일어났고, 다음에 무엇이 일어날지'를 예측하고 제안하는 미래 지향적 분석(예측적/처방적 분석)에 중점을 둡니다.
\end{summarybox}

\begin{table}[h!]
\centering
\caption{BI와 BA의 비교}
\label{tab:bi_ba_comparison}
\begin{adjustbox}{width=\textwidth,center}
\begin{tabular}{p{0.15\textwidth}p{0.4\textwidth}p{0.4\textwidth}}
\toprule
\textbf{특징} & \textbf{비즈니스 인텔리전스 (BI)} & \textbf{비즈니스 분석 (BA)} \\
\midrule
\textbf{주요 질문} & 무엇이 일어났는가? 언제? 누가? 얼마나 많이? 어떻게? & 왜 일어났는가? 다시 일어날까? X를 바꾸면 어떻게 될까? 다음에 무엇을 해야 할까? \\
\textbf{분석 유형} & 기술적 분석 (Descriptive Analysis) & 예측적 분석 (Predictive Analysis), 처방적 분석 (Prescriptive Analysis) \\
\textbf{데이터 사용} & 과거 및 현재 데이터 & 과거 데이터를 사용하여 현재를 설명하고 미래를 예측 \\
\textbf{목표} & 현재 비즈니스 상황 이해, 보고서, 대시보드 & 미래 행동 제안, 의사결정 지원, 최적화 \\
\textbf{예시} & 월별 매출 보고서, 웹사이트 트래픽 현황 & 다음 분기 매출 예측, 특정 마케팅 캠페인의 효과 예측 \\
\bottomrule
\end{tabular}
\end{adjustbox}
\end{table}

\subsubsection{BI의 현대적 트렌드: 처방적 분석으로의 전환}
전통적인 BI는 주로 기술적 분석(Descriptive Analysis)에 중점을 두었지만, 현재는 처방적 분석(Prescriptive Analysis)으로 나아가는 추세입니다.

\begin{itemize}
    \item \textbf{처방적 분석 (Prescriptive Analysis):} 단순히 '무엇이 일어났는지'를 넘어, '무엇을 해야 하는지'에 대한 구체적인 제안을 제공하여 비즈니스 이해관계자의 의사결정을 돕습니다.
    \item \textbf{주요 요소:}
    \begin{itemize}
        \item \textbf{실행 가능한 통찰력 (Actionable Insights):} 비즈니스 이해관계자가 즉시 행동할 수 있는 구체적인 제안.
        \item \textbf{셀프 서비스 모니터링 (Self-service Monitoring):} 비즈니스 사용자가 직접 도구 및 서비스 사용량을 모니터링.
        \item \textbf{성능 최적화 (Performance Optimization):} 도구 개선을 위한 지속적인 노력.
        \item \textbf{보안 제어 (Security Controls):} 데이터 접근 권한 관리 및 보안 유지.
        \item \textbf{스토리텔링 (Storytelling):} 비즈니스 사용자가 정보를 쉽게 이해하고 소화할 수 있도록 이야기 형식으로 공유.
    \end{itemize}
\end{itemize}

\subsection{BI 프로세스 단계}
BI는 데이터를 수집하는 것부터 시작하여 최종적으로 비즈니스 의사결정에 활용되기까지 여러 단계를 거칩니다.

\begin{summarybox}[title=핵심 요약: 데이터에서 의사결정까지]
BI 프로세스는 데이터를 모으고(수집), 정리하고(조직), 분석하고(분석), 시각화하여(시각화), 최종적으로 비즈니스 리더들이 중요한 결정을 내릴 수 있도록(공유 및 활용) 돕는 일련의 체계적인 과정입니다.
\end{summarybox}

\begin{enumerate}
    \item \textbf{데이터 수집 (Collecting Data):}
    \begin{itemize}
        \item 조직 내 여러 소스에서 데이터를 수집하여 데이터 웨어하우스와 같은 중앙 저장소로 가져옵니다.
    \end{itemize}
    \item \textbf{데이터 조직 (Organizing Data):}
    \begin{itemize}
        \item 수집된 데이터를 차원 모델과 같은 데이터 모델로 구성합니다.
    \end{itemize}
    \item \textbf{데이터 분석 (Analyzing Data):}
    \begin{itemize}
        \item SQL 쿼리 또는 다른 유형의 쿼리를 사용하여 데이터를 분석하고 시각화 및 대시보드를 구축합니다.
        \item 데이터 마이닝, 예측 분석, 통계 분석, 대규모 데이터 시각화, KPI(핵심 성과 지표), 성능 벤치마킹, 임시 쿼리 등이 포함됩니다.
    \end{itemize}
    \item \textbf{결과 공유 (Sharing Results):}
    \begin{itemize}
        \item 분석 결과를 비즈니스 이해관계자에게 공유합니다. 대시보드, 스코어카드, 임시 쿼리 등을 통해 공유될 수 있습니다.
    \end{itemize}
    \item \textbf{비즈니스 활용 (Business Utilization):}
    \begin{itemize}
        \item 비즈니스 임원 및 기타 관계자들이 이 데이터를 사용하여 제품 개발, 시장 진출, 제품 취소 등과 같은 의사결정을 내립니다.
    \end{itemize}
\end{enumerate}

\subsubsection{좋은 데이터 관리의 구성 요소}
\begin{itemize}
    \item \textbf{데이터 수집 (Data Collection):} 여러 소스에서 데이터 수집.
    \item \textbf{데이터 준비 (Data Preparation):} 통합, 모델링, 클렌징, 보강.
    \item \textbf{분석 (Analytics):} '무엇이 일어났고, 왜 일어났으며, 어떻게 일어났는지, 미래에 무엇이 일어날지' 탐색.
    \item \textbf{데이터 시각화 (Data Visualization):} 차트 및 그래프 제공.
    \item \textbf{공유 및 협업 (Sharing and Collaboration):} 정보 공유.
    \item \textbf{거버넌스 (Governance):} 데이터 접근 권한 관리.
    \item \textbf{전략 문서화 (Strategy Documentation):} 비즈니스 목표 및 전략 문서화.
    \textbf{중앙 집중식 데이터 카탈로그 (Centralized Data Catalogs):} 데이터 위치 및 메타데이터 관리.
    \item \textbf{사용 편의성 (Ease of Use):} 비즈니스 사용자의 데이터 접근 용이성.
\end{itemize}

\subsection{BI 분석가 페르소나}
데이터 엔지니어링 분야에서는 다양한 페르소나가 존재하며, BI 분석가는 그중 하나로 특정 기술 세트와 역할을 가집니다.

\begin{summarybox}[title=핵심 요약: 비즈니스 통찰력의 전달자]
BI 분석가는 주로 SQL을 사용하여 골드(Gold) 또는 실버(Silver) 계층의 정제된 데이터를 분석하고, 대시보드와 보고서를 통해 비즈니스 사용자에게 실행 가능한 통찰력을 제공하는 역할을 합니다. 이들은 데이터 민주화의 핵심 주체입니다.
\end{summarybox}

\begin{itemize}
    \item \textbf{주요 기술 세트:} SQL 사용 능력.
    \item \textbf{사용 데이터:} 주로 큐레이션된 데이터(골드 계층 또는 실버/골드 계층).
    \item \textbf{데이터 저장소:} 전통적으로 데이터 웨어하우스 및 특수 데이터 마트(Data Marts)에 저장된 데이터를 사용합니다.
    \item \textbf{레이크하우스의 역할:} 이제 레이크하우스는 데이터 레이크와 웨어하우스의 역할을 모두 수행하여 BI 분석가에게 필요한 데이터를 직접 지원할 수 있습니다.
    \item \textbf{데이터 민주화 (Democratization of Data):} 셀프 서비스 BI를 통해 개인이 직접 데이터를 탐색하고, '만약(What If)' 분석을 위한 데이터 마이닝을 수행할 수 있도록 돕습니다.
\end{itemize}

\subsection{데이터 웨어하우스 용어}
데이터 웨어하우스 환경에서 사용되는 주요 용어들을 이해하는 것은 BI 분석에 필수적입니다.

\begin{summarybox}[title=핵심 요약: 데이터 구조의 언어]
데이터 웨어하우스 용어는 데이터가 어떻게 조직되고, 저장되며, 접근되고, 관계를 맺는지 설명하는 '언어'입니다. 이 용어들을 이해하면 데이터의 구조와 품질을 파악하고 효율적으로 쿼리할 수 있습니다.
\end{summarybox}

\begin{itemize}
    \item \textbf{카탈로그 (Catalog):} 메타데이터(데이터에 대한 데이터)를 저장합니다.
    \item \textbf{데이터베이스/스키마 (Database/Schema):} 네임스페이스를 제공하고 데이터의 논리적 및 물리적 조직을 정의합니다.
    \item \textbf{테이블 (Tables):}
    \begin{itemize}
        \item \textbf{관리형 테이블 (Managed Tables):} 데이터 레이크에 의해 관리됩니다.
        \item \textbf{비관리형 테이블 (Unmanaged Tables):} 레이크 외부에 저장되며 메타데이터를 통해 접근됩니다.
    \end{itemize}
    \item \textbf{키 (Keys):}
    \begin{itemize}
        \item \textbf{기본 키 (Primary Keys):} 데이터 내 엔티티를 고유하게 식별합니다.
        \item \textbf{외래 키 (Foreign Keys):} 데이터 간의 관계를 설정합니다.
        \textbf{대리 키 (Surrogate Keys):} 비즈니스 의미가 없는 인공적인 키.
        \item \textbf{자연 키 (Natural Keys):} 비즈니스 의미를 가지는 키.
    \end{itemize}
    \item \textbf{제약 조건 (Constraints):} 데이터 품질을 보장하기 위해 데이터에 적용되는 규칙입니다.
    \item \textbf{뷰 (Views):}
    \begin{itemize}
        \item 가상 테이블로, 접근 시점에 기본 쿼리가 실행되어 데이터가 생성됩니다. 저장 공간을 차지하지 않습니다.
    \end{itemize}
    \item \textbf{페더레이션 쿼리 (Federated Queries):}
    \begin{itemize}
        \item 여러 데이터 저장소 경계를 넘어 데이터를 요청할 수 있으며, '푸시다운 프레디케이트(Pushdown Predicates)'를 통해 효율적으로 데이터를 가져옵니다.
    \end{itemize}
    \item \textbf{구체화된 뷰 (Materialized Views):}
    \begin{itemize}
        \item 뷰와 유사하지만 미리 계산되어 저장됩니다. 데이터에 더 빠른 접근을 지원합니다.
    \end{itemize}
    \item \textbf{저장 프로시저/UDF (Stored Procedures/UDFs):}
    \begin{itemize}
        \item 특정 로직을 캡슐화하여 재사용 가능하게 합니다.
    \end{itemize}
    \item \textbf{시맨틱 데이터 모델 (Semantic Data Model):}
    \begin{itemize}
        \item 비즈니스 엔티티 간의 관계를 비즈니스 용어(예: 고객, 제품, 거래)를 사용하여 캡처합니다.
    \end{itemize}
\end{itemize}

\subsection{데이터 모델링 전략}
데이터 모델링은 데이터를 구조화하여 효율적인 저장, 검색 및 분석을 가능하게 하는 과정입니다. BI 및 분석 엔지니어링에서 중요한 역할을 합니다.

\begin{summarybox}[title=핵심 요약: 데이터의 청사진]
데이터 모델링은 데이터의 '청사진'을 그리는 과정으로, 데이터를 어떻게 저장하고 연결할지 결정합니다. 이는 데이터의 일관성을 유지하고, 쿼리 성능을 최적화하며, 비즈니스 사용자가 데이터를 쉽게 이해하고 활용할 수 있도록 돕습니다.
\end{summarybox}

\begin{itemize}
    \item \textbf{객체 지향 프로그래밍의 개념:} 'is a(상속)' 및 'has a(집합/연관)' 관계를 통해 엔티티와 관계를 파악합니다. 명사(엔티티)와 동사(관계)를 식별합니다.
    \item \textbf{모델링 단계:}
    \begin{enumerate}
        \item \textbf{시맨틱 모델 (Semantic Model):} 비즈니스 관점에서 '무엇을 해야 하는지' 이해합니다.
        \item \textbf{논리적 모델 (Logical Model):} 시맨틱 모델을 기술적인 논리적 구조로 변환합니다.
        \item \textbf{물리적 모델 (Physical Model):} 특정 기술 스택(데이터베이스, 플랫폼)에 맞춰 논리적 모델을 구현합니다.
    \end{enumerate}
    \item \textbf{주요 패턴:}
    \begin{itemize}
        \item \textbf{제3 정규형 (Third Normal Form, 3NF) - Bill Inmon 전략:}
        \begin{itemize}
            \item 데이터 중복을 최소화하고 데이터 무결성을 높이는 데 중점을 둡니다.
            \item 엄격한 참조 무결성(Referential Integrity)을 유지합니다.
            \item 주로 OLTP(온라인 트랜잭션 처리) 시스템에 적합하지만, 분석 쿼리에는 복잡한 조인이 필요할 수 있습니다.
        \end{itemize}
        \item \textbf{차원 모델링 (Dimensional Modeling) - Ralph Kimball 전략:}
        \begin{itemize}
            \item BI 및 분석 사용 사례에 최적화된 모델링 방식입니다.
            \item \textbf{스타 스키마 (Star Schema):} 중앙에 팩트(Fact) 테이블이 있고, 그 주위에 차원(Dimension) 테이블이 직접 연결된 형태입니다. 쿼리 성능이 빠르고 이해하기 쉽습니다.
            \item \textbf{스노우플레이크 스키마 (Snowflake Schema):} 스타 스키마에서 차원 테이블이 추가적으로 정규화되어 더 많은 차원 테이블로 분리된 형태입니다. 데이터 중복을 줄이지만, 쿼리가 더 복잡해질 수 있습니다.
            \item 쿼리는 이러한 BI 사용 사례를 지원하도록 최적화됩니다.
        \end{itemize}
        \item \textbf{데이터 볼트 (Data Vault):}
        \begin{itemize}
            \item 더 유연하고 쓰기(Write) 작업을 허용하도록 설계되었습니다.
            \item \textbf{허브 (Hubs):} 제품, 고객과 같은 핵심 비즈니스 개념을 저장합니다.
            \textbf{링크 (Links):} 다양한 허브 간의 관계(기본 키와 외래 키)를 정의합니다.
            \item \textbf{새틀라이트 (Satellites):} 비즈니스 시맨틱스(설명 속성)를 저장합니다.
            \item 실버(Silver) 계층에 데이터 볼트를 사용할 수 있지만, 골드(Gold) 계층의 보고서 레이어는 일반적으로 차원 모델로 다시 변환됩니다.
        \end{itemize}
    \end{itemize}
\end{itemize}

\subsection{데이터브릭스 플랫폼 아키텍처}
Databricks 플랫폼은 데이터 엔지니어링, 데이터 과학, BI/분석을 위한 통합된 환경을 제공합니다.

\begin{summarybox}[title=핵심 요약: 통합된 데이터 여정]
Databricks 플랫폼은 데이터 수집부터 변환, 쿼리, 분석, 거버넌스, 공유에 이르는 데이터의 전체 여정을 단일 플랫폼에서 지원합니다. 특히 Databricks SQL은 BI 분석가에게 친숙한 SQL 환경에서 고성능 분석을 가능하게 합니다.
\end{summarybox}

\begin{figure}[h!]
    \centering
    % \includegraphics[width=0.9\textwidth]{databricks_architecture_conceptual.png}  % Image not found: databricks_architecture_conceptual.png % 이미지 파일이 없으므로 개념적 설명으로 대체
    \caption{Databricks 플랫폼의 개념적 아키텍처 (이미지 없음)}
    \label{fig:databricks_architecture}
\end{figure}

\begin{itemize}
    \item \textbf{소스 (Sources):}
    \begin{itemize}
        \item 정형, 반정형, 비정형 데이터가 여러 클라우드에서 유입됩니다.
        \item \textbf{ETL (Extract, Transform, Load):} 데이터 레이크가 데이터의 소유자가 되며, Bronze, Silver, Gold 변환을 거칩니다.
        \item \textbf{페더레이션 (Federation):} 읽기 전용 데이터를 다른 저장소(예: 온프레미스 Oracle, Amazon Redshift)에서 가져옵니다. 데이터 소유권은 원본에 유지되며, 물리적으로 레이크로 가져오지 않습니다. 주로 참조 데이터나 소규모 데이터를 통합하는 데 사용됩니다. 이는 '데이터 가상화'의 한 예시입니다.
    \end{itemize}
    \item \textbf{수집 (Ingest):} 배치 및 스트리밍 데이터 수집 (Autoloader, DLT, Copy Into 등).
    \item \textbf{변환 (Transformation):} 데이터 클렌징, 조인, 정규화, 집계 등 핵심 데이터 처리 (Bronze, Silver, Gold 계층).
    \item \textbf{쿼리 및 처리 (Query and Processing):}
    \begin{itemize}
        \item \textbf{Databricks SQL (DBSQL):} 데이터 웨어하우징을 위한 부분으로, 표준 SQL을 사용하여 데이터를 쿼리하고 변환합니다. LLM을 SQL 내에서 사용하는 강력한 기능도 제공합니다.
        \item \textbf{ML (Machine Learning):} 데이터에 대한 ML 모델 실행 및 통찰력(insights)을 DBSQL 대시보드를 통해 시각화합니다.
    \end{itemize}
    \item \textbf{스토리지 (Storage):}
    \begin{itemize}
        \item 데이터는 항상 클라우드에 저장됩니다 (Databricks는 멀티클라우드 지원).
        \item Delta Lake 프로토콜은 ACID(원자성, 일관성, 고립성, 지속성) 준수, 최적화, 신뢰성을 보장합니다.
        \item Unity Catalog는 스키마, 설명, 계보 정보와 같은 메타데이터를 관리합니다.
    \end{itemize}
    \item \textbf{협업 (Collaboration):}
    \begin{itemize}
        \item \textbf{Delta Sharing:} Databricks 플랫폼을 사용하지 않는 외부 사용자에게 데이터를 공유하는 방법.
        \item \textbf{Marketplace:} 타사 데이터(날씨, 신용 정보 등)를 가져오는 방법.
    \end{itemize}
    \item \textbf{거버넌스 (Governance):}
    \begin{itemize}
        \item \textbf{Unity Catalog:} 데이터, 모델, 피처에 대한 카탈로그화, 데이터 계보(Lineage), 접근 제어를 담당하는 통합 거버넌스 레이어입니다.
        \item \textbf{메타스토어 (Metastore):} Unity Catalog의 계층 구조에서 최상위 논리적 구성 요소로, 카탈로그, 스키마, 테이블 등을 포함하는 컨테이너 역할을 합니다.
    \end{itemize}
    \item \textbf{인텔리전스 엔진 (Intelligence Engine):}
    \begin{itemize}
        \item \textbf{Databricks IQ:} 플랫폼 사용을 관찰하고 최적화를 돕는 지능형 엔진입니다.
        \item \textbf{Assistant:} 노트북 및 SQL 편집기에서 코드 생성 및 질문 답변을 돕습니다.
        \item \textbf{Genie Rooms:} 구조화된 데이터에 대해 자연어로 질문하고 SQL 생성, 시각화, 설명 등을 제공합니다.
    \end{itemize}
    \item \textbf{컴퓨팅 (Compute):}
    \begin{itemize}
        \item \textbf{Spark:} 데이터 처리의 핵심 엔진.
        \item \textbf{Photon:} C++로 재작성된 Spark의 차세대 버전으로, 벡터화 및 고성능을 제공합니다. DBSQL 및 서버리스 환경에서 자동으로 사용됩니다.
        \item \textbf{DLT (Delta Live Tables) / Lakeflow:} Spark 선언적 파이프라인으로, 데이터 파이프라인 구축을 단순화합니다.
    \end{itemize}
    \item \textbf{분석 (Analysis):}
    \begin{itemize}
        \item 대시보드, 레이크하우스 앱, 외부 BI 도구(Tableau, PowerBI)를 통한 데이터 시각화 및 분석.
    \end{itemize}
    \item \textbf{통합 (Integrations):}
    \begin{itemize}
        \item ID 공급자, 엔터프라이즈 카탈로그(Alation, Collibra), 클라우드 네이티브 서비스(AWS Bedrock), Hugging Face 등 다양한 외부 시스템과의 통합.
    \end{itemize}
\end{itemize}

\begin{examplebox}[title=페더레이션과 ETL의 차이]
\textbf{ETL:} 마치 이사하는 것과 같습니다. 데이터를 원본 집에서 새 집(데이터 레이크)으로 완전히 옮겨와서 소유권을 가지고 마음대로 정리(변환)하고 사용합니다.
\textbf{페더레이션:} 마치 도서관에서 책을 빌려보는 것과 같습니다. 책(데이터)은 여전히 도서관(원본 시스템)에 있고, 나는 그 책을 잠시 빌려와서 내 노트(레이크의 다른 데이터)와 함께 참고할 뿐, 책의 소유권을 가지거나 마음대로 내용을 바꿀 수 없습니다. 주로 소규모 참조 데이터에 사용됩니다.
\end{examplebox}

\subsection{Databricks SQL의 주요 기능}
Databricks SQL은 Databricks 레이크하우스 플랫폼의 핵심 구성 요소로, BI 및 데이터 웨어하우징 워크로드를 위해 설계되었습니다.

\begin{summarybox}[title=핵심 요약: 레이크하우스의 BI 엔진]
Databricks SQL은 레이크하우스의 데이터를 직접 활용하여 고성능 BI 분석을 가능하게 하는 강력한 SQL 기반 환경입니다. 최신 데이터에 대한 대시보드, AI 기반 분석, 그리고 다양한 외부 BI 도구와의 통합을 지원하여 데이터 민주화를 실현합니다.
\end{summarybox}

\begin{itemize}
    \item \textbf{최신 데이터 분석:} 파이프라인에 의해 최신 데이터가 수집되면, 워크플로우의 한 노드로 대시보드를 생성하여 자동으로 대시보드를 새로 고칠 수 있습니다.
    \item \textbf{BI 분석가 친화적:} 데이터 엔지니어뿐만 아니라 데이터 과학자 및 BI 분석가를 위한 홈 역할을 합니다. 표준 SQL(NCSQL)을 사용하여 코드가 매우 이식성이 좋습니다.
    \item \textbf{개방형 기반:} 개방형 표준을 기반으로 하여 유연성을 제공합니다.
    \item \textbf{광범위한 BI 도구 통합:} Tableau, PowerBI, Looker와 같은 외부 BI 도구가 레이크 데이터에 직접 연결하여 최신 데이터에 대한 정교한 시각화를 수행할 수 있습니다.
    \item \textbf{가격 대비 성능:} 다른 클라우드 데이터 웨어하우스보다 우수한 가격 대비 성능을 제공합니다.
    \item \textbf{간소화된 발견 및 공유:} Unity Catalog를 통해 새로운 통찰력의 발견 및 공유를 간소화합니다. 통합 검색 인터페이스, 간소화된 관리 및 거버넌스, 데이터 계보 보기를 제공합니다.
    \item \textbf{Lakeview 대시보드 및 Genie Spaces:} 대시보드 및 자연어 기반의 Genie 기능을 지원합니다.
    \item \textbf{지능형 데이터 웨어하우징:} 자연어로 데이터에 질문하고 답변을 얻을 수 있는 기능을 제공합니다.
    \item \textbf{자동화된 관리 및 튜닝:} 서버리스(Serverless) 사용으로 백그라운드에서 최적화가 자동으로 이루어지며, 최적의 총 소유 비용(TCO)을 제공합니다.
    \item \textbf{탐색적 데이터 분석 (EDA):} SQL 함수를 사용하여 편집기에서 지능형 자동 완성 및 NCSQL을 통해 EDA를 수행할 수 있습니다.
    \item \textbf{쿼리 기록 및 프로파일링:} 쿼리 기록을 확인하고 쿼리 프로파일을 분석하여 성능 최적화를 수행할 수 있습니다.
    \item \textbf{연결성:} JDBC/ODBC, Python, Java, NodeJS, Go 등 다양한 언어 및 도구와의 연결을 지원합니다.
    \textbf{성능:} Photon 엔진과 예측 IO(Predictive IO)를 통해 고성능을 제공합니다.
    \item \textbf{쿼리 페더레이션 및 구체화된 뷰:} 여러 소스에 걸친 조인 및 미리 계산된 뷰를 통해 데이터 접근을 가속화합니다.
    \item \textbf{워크플로우 통합:} SQL 태스크를 대시보드에 직접 연결하여 데이터 파이프라인과 대시보드를 통합합니다.
    \textbf{Python UDFs:} SQL 내에서 Python UDF를 사용할 수 있습니다.
    \item \textbf{지리공간 기능 (Geospatial Capabilities):} H3 지원 및 SQL 내 지리공간 함수를 통해 공간 데이터 분석을 지원합니다.
    \item \textbf{서버리스 데이터 웨어하우스:} VM(가상 머신)이 미리 준비되어 있어 쿼리 요청 시 즉시 컴퓨팅을 제공하며, 높은 동시성(High Concurrency) BI 워크로드를 처리하고 결과 캐싱을 지원합니다.
    \item \textbf{거버넌스 및 보안:} Unity Catalog에 의해 거버넌스 및 보안이 제공됩니다. 행 수준 보안(Row-level Security) 및 열 마스킹(Column Masking)을 통해 세분화된 접근 제어를 지원합니다.
    \item \textbf{데이터 검색:} 통합 검색을 통해 노트북, 대시보드 등 다양한 데이터 자산을 검색할 수 있습니다.
    \item \textbf{감사 (Auditing):} 누가 어떤 테이블에 접근했는지, 성공 여부, 사용 시간 등을 기록하여 규제 산업의 요구사항을 충족하고 내부 위협을 감지하는 데 사용됩니다.
    \item \textbf{스트리밍 테이블 (Streaming Tables):} 실시간 사용 사례를 위해 클라우드 스토리지나 메시지 큐(Kafka, Kinesis)에서 들어오는 새로운 데이터를 SQL을 통해 스트리밍 테이블로 생성할 수 있습니다.
    \item \textbf{AI 쿼리 및 AI 함수:} SQL 쿼리 내에서 LLM 모델을 호출하여 데이터에 대한 감성 분석, 분류, 정보 추출, 문법 수정, 데이터 마스킹, 유사성 분석 등의 작업을 수행할 수 있습니다.
\end{itemize}

\subsection{BI의 핵심 KPI: 동시성 및 확장성}
BI 사용 사례에서 플랫폼의 성능을 평가하는 두 가지 중요한 지표는 동시성(Concurrency)과 확장성(Scaling)입니다.

\begin{summarybox}[title=핵심 요약: BI 시스템의 체력과 유연성]
동시성은 시스템이 동시에 얼마나 많은 요청을 지연 없이 처리할 수 있는지를 나타내고, 확장성은 데이터 양이나 사용자 수가 늘어날 때 시스템이 얼마나 유연하게 성능을 유지할 수 있는지를 나타냅니다. 이 두 가지는 BI 대시보드의 원활한 작동에 필수적입니다.
\end{summarybox}

\begin{itemize}
    \item \textbf{동시성 (Concurrency):}
    \begin{itemize}
        \item \textbf{정의:} 시스템이 느려지거나 응답하지 않게 되지 않고 동시에 처리할 수 있는 쿼리(요청)의 수입니다.
        \item \textbf{중요성:} 대시보드에는 여러 위젯이 있으며, 이들은 각각 다른 쿼리에 의해 데이터가 채워집니다. 많은 사용자가 동시에 대시보드에 접근할 때, 시스템이 모든 쿼리를 효율적으로 처리해야 합니다.
        \item \textbf{예시:} 200명의 분석가가 동시에 SQL 엔드포인트에 쿼리를 날릴 때, 각 쿼리가 대기 시간 없이 즉시 실행되어야 합니다.
    \end{itemize}
    \item \textbf{확장성 (Scaling):}
    \begin{itemize}
        \item \textbf{정의:} 데이터 양이 증가하거나 사용자 수가 늘어날 때, 시스템이 그에 맞춰 성능을 유지하거나 향상시킬 수 있는 능력입니다.
        \item \textbf{중요성:} 데이터는 계속해서 증가하므로, BI 플랫폼은 이러한 데이터 증가에 맞춰 컴퓨팅 자원을 유연하게 확장할 수 있어야 합니다.
        \item \textbf{예시:} 5GB의 데이터에서 20GB로 데이터가 늘어나도, 대시보드 로딩 속도가 크게 저하되지 않아야 합니다.
    \end{itemize}
\end{itemize}

\begin{warningbox}[title=쿼리 프로파일의 중요성]
Databricks SQL의 쿼리 프로파일에서 쿼리가 실행되는 시간(녹색)과 대기하는 시간(노란색)을 확인할 수 있습니다. 이상적인 시스템은 대부분 녹색이고 노란색은 거의 없어야 합니다. 이는 컴퓨팅 자원이 낭비되지 않고 쿼리가 제때 처리되고 있음을 의미합니다.
\end{warningbox}

\newpage
\section{절차/방법}

\subsection{데이터 레이크 하이드레이션 방법}
데이터 레이크에 데이터를 채우는(Hydrating) 다양한 방법들이 있습니다. 이는 데이터의 유형, 속도, 양에 따라 적절한 방법을 선택해야 합니다.

\begin{summarybox}[title=핵심 요약: 데이터 레이크에 물 채우기]
데이터 레이크에 데이터를 채우는 것은 마치 저수지에 물을 채우는 것과 같습니다. REST API, 스트리밍, 자동 로더, COPY INTO, MERGE/UPSERT 등 다양한 방법으로 데이터를 가져올 수 있으며, 각 방법은 데이터의 특성과 요구사항에 맞춰 선택됩니다.
\end{summarybox}

\begin{enumerate}
    \item \textbf{REST API:}
    \begin{itemize}
        \item 웹 서비스를 통해 데이터를 가져오는 일반적인 방법입니다.
    \end{itemize}
    \item \textbf{스트리밍 데이터 (Streaming Data):}
    \begin{itemize}
        \item 실시간 또는 거의 실시간으로 지속적으로 발생하는 데이터를 처리하는 방법입니다.
        \item 예: Kafka, Kinesis와 같은 메시지 큐 시스템에서 데이터를 가져옵니다.
    \end{itemize}
    \item \textbf{자동 로더 (Autoloader):}
    \begin{itemize}
        \item 클라우드 스토리지(예: S3, ADLS)에 새로운 파일이 도착하는 것을 자동으로 감지하고 읽어 데이터 레이크로 가져오는 기능입니다.
        \item 데이터 파이프라인을 자동화하고 관리 오버헤드를 줄이는 데 매우 유용합니다.
    \end{itemize}
    \item \textbf{COPY INTO:}
    \begin{itemize}
        \item 클라우드 스토리지에서 Delta Lake 테이블로 데이터를 효율적으로 로드하는 SQL 명령입니다.
        \item 대규모 배치 데이터 로드에 적합하며, 중복 방지 기능을 제공합니다.
    \end{itemize}
    \item \textbf{MERGE/UPSERT (병합/업서트):}
    \begin{itemize}
        \item 기존 테이블에 새로운 데이터를 삽입(INSERT)하거나 기존 데이터를 업데이트(UPDATE)하는 작업을 단일 명령으로 수행하는 방법입니다.
        \item 변경 데이터 캡처(Change Data Capture, CDC) 시나리오에서 특히 유용하며, 데이터의 일관성을 유지하는 데 도움을 줍니다.
    \end{itemize}
\end{enumerate}

\subsection{Databricks SQL 대시보드 생성 및 활용}
Databricks SQL은 대시보드를 생성하고 관리하며, 이를 통해 비즈니스 통찰력을 시각화하고 공유하는 강력한 기능을 제공합니다.

\begin{summarybox}[title=핵심 요약: 대시보드로 데이터 스토리텔링]
Databricks SQL 대시보드는 데이터를 시각적으로 표현하고, 필터와 텍스트 박스를 통해 맥락을 제공하며, 심지어 AI 어시스턴트(Genie)를 활용하여 자연어로 질문하고 답변을 얻을 수 있는 동적인 도구입니다. 이를 통해 비즈니스 사용자는 데이터를 쉽게 이해하고 의사결정에 활용할 수 있습니다.
\end{summarybox}

\begin{enumerate}
    \item \textbf{대시보드 접근:} Databricks 워크스페이스에서 'Dashboards' 섹션으로 이동합니다.
    \item \textbf{기존 대시보드 탐색:}
    \begin{itemize}
        \item 'New York City taxi trip analysis' 및 'Workspace usage'와 같은 기본 제공 대시보드를 확인할 수 있습니다.
        \item 대시보드에는 'Data' 탭과 'Summary' 탭이 있습니다. 'Data' 탭에서는 대시보드를 채우는 데 사용되는 데이터셋(예: trips, route revenue)을 볼 수 있습니다.
    \end{itemize}
    \item \textbf{대시보드 구성 요소:}
    \begin{itemize}
        \item \textbf{시각화 위젯 (Visualization Widget):} 차트, 그래프 등 데이터를 시각적으로 표현합니다.
        \item \textbf{텍스트 박스 (Text Box):} 대시보드에 설명, 제목, 추가 정보를 제공합니다. 크기, 색상, 글꼴 등을 편집할 수 있습니다.
        \item \textbf{필터 (Filter):} 대시보드 데이터를 동적으로 필터링할 수 있는 드롭다운 메뉴 등을 제공합니다. 특정 기간, 지역, 카테고리 등으로 데이터를 제한할 수 있습니다.
    \end{itemize}
    \item \textbf{새로운 시각화 생성 (AI 어시스턴트 활용):}
    \begin{itemize}
        \item 'Add a data visualization' 버튼을 클릭합니다.
        \item AI 어시스턴트가 자동으로 '총 수익', '경로별 수익 분포', '픽업 지역별 여행 횟수' 등 다양한 시각화 아이디어를 제안합니다.
        \item 원하는 제안을 선택하면, 어시스턴트가 자동으로 시각화(예: 막대 차트)를 생성합니다.
        \item 생성된 시각화의 제목, 설명, 데이터셋, 차트 유형(막대, 산점도 등), 색상 등을 왼쪽 패널에서 수정할 수 있습니다.
        \item 어시스턴트가 생성한 SQL 코드를 확인하고, 필요에 따라 추가적인 지시를 내리거나 수락/거부할 수 있습니다.
    \end{itemize}
    \item \textbf{대시보드 게시 (Publish):}
    \begin{itemize}
        \item 대시보드를 '개발자 뷰'에서 '게시된 뷰'로 전환할 수 있습니다.
        \item \textbf{내부 공유:} Databricks 워크스페이스 내 다른 사용자에게 '관리' 또는 '보기' 권한으로 공유할 수 있습니다.
        \item \textbf{외부 공유:} '임베디드 자격 증명(Embedded Credentials)'을 사용하여 Databricks 플랫폼에 온보딩되지 않은 외부 사용자에게 대시보드를 게시할 수 있습니다. 이 경우, 대시보드 새로 고침 시 발생하는 컴퓨팅 비용은 게시자가 부담합니다.
    \end{itemize}
    \item \textbf{새로 고침 및 경고 설정:}
    \begin{itemize}
        \item 대시보드 새로 고침 일정을 설정할 수 있습니다 (예: 15분마다).
        \item 워크플로우의 태스크로 대시보드 새로 고침을 포함하여, 새로운 데이터가 들어올 때 자동으로 대시보드가 업데이트되도록 할 수 있습니다.
        \item 특정 쿼리 결과가 임계값(예: 행 수가 특정 값 미만)을 벗어날 경우, Slack, PagerDuty, 이메일 등으로 경고를 받을 수 있습니다. 이는 데이터 파이프라인의 문제나 비정상적인 상황을 조기에 감지하는 데 유용합니다.
    \end{itemize}
\end{enumerate}

\subsection{Genie 활용법}
Genie는 Databricks 대시보드에 통합된 AI 기능으로, 비즈니스 사용자가 자연어로 데이터에 질문하고 통찰력을 얻을 수 있도록 돕습니다.

\begin{summarybox}[title=핵심 요약: 자연어로 데이터와 대화하기]
Genie는 비즈니스 사용자가 SQL을 몰라도 대시보드 데이터에 대해 자연어로 질문하고, Genie가 자동으로 SQL 쿼리를 생성하여 답변과 시각화를 제공하는 '데이터 대화 도우미'입니다. 이는 비즈니스 사용자의 데이터 탐색 자율성을 크게 높여줍니다.
\end{summarybox}

\begin{enumerate}
    \item \textbf{Genie 활성화:} 게시된 대시보드 뷰에서 'Ask Genie some questions' 아이콘을 클릭하여 Genie를 활성화합니다.
    \item \textbf{자연어 질문:} Genie 채팅 인터페이스에 자연어로 질문을 입력합니다 (예: "What is an average trip duration?").
    \item \textbf{Genie의 답변:} Genie는 질문을 분석하여 답변을 제공합니다 (예: "평균 여행 시간은 13.7분입니다.").
    \item \textbf{SQL 코드 확인:} Genie가 답변을 생성하기 위해 사용한 SQL 쿼리를 확인할 수 있습니다. 이를 통해 Genie의 답변이 어떻게 도출되었는지 투명하게 이해할 수 있습니다.
    \item \textbf{지시 세트 (Instruction Sets) 활용:}
    \begin{itemize}
        \item Genie가 질문을 정확하게 이해하지 못하거나 특정 방식으로 답변하기를 원할 경우, '지시 세트'를 구성할 수 있습니다.
        \item 예를 들어, "여행 시간은 분이 아닌 초 단위로 계산해줘"와 같은 추가 지시를 제공할 수 있습니다.
        \item 특정 복잡한 쿼리에 대해 Genie가 잘못된 답변을 하지 않도록, 미리 정의된 SQL 함수를 사용하도록 지시할 수도 있습니다.
    \end{itemize}
\end{enumerate}

\begin{examplebox}[title=Genie의 강점]
Genie는 Chat GPT와 달리, 인터넷의 일반적인 지식이 아닌 \textbf{주어진 테이블의 메타데이터와 데이터셋만을 기반으로 답변}합니다. 따라서 '환각(Hallucination)' 현상이 거의 없으며, 질문이 불분명할 경우 답변을 거부하거나 명확화를 요청합니다. 이는 비즈니스 사용자가 신뢰할 수 있는 데이터 기반 통찰력을 얻는 데 매우 중요합니다.
\end{examplebox}

\newpage
\section{실습/코드/오류와 해결}

강의에서는 Databricks 플랫폼을 활용한 실제 대시보드 및 SQL 편집기 사용 시연이 이루어졌습니다. 여기서는 시연에서 다루어진 주요 기능과 쿼리 프로파일 분석 방법을 설명합니다. 실제 코드는 제공되지 않았으므로, 기능 설명에 중점을 둡니다.

\subsection{Databricks SQL 편집기 활용}
Databricks SQL 편집기는 BI 분석가가 쿼리를 작성하고 실행하며, 데이터를 탐색하는 데 사용되는 핵심 도구입니다.

\begin{summarybox}[title=핵심 요약: SQL 분석가의 작업대]
Databricks SQL 편집기는 단순한 텍스트 편집기를 넘어, 자동 완성, 멀티탭 기능, AI 어시스턴트 통합, 그리고 SQL 웨어하우스(컴퓨팅) 연결을 통해 효율적인 SQL 쿼리 작성 및 실행 환경을 제공합니다.
\end{summarybox}

\begin{itemize}
    \item \textbf{기능:}
    \begin{itemize}
        \item \textbf{자동 완성 (Autocomplete):} SQL 키워드, 테이블, 컬럼 이름 등을 자동으로 제안하여 쿼리 작성 속도를 높입니다.
        \item \textbf{멀티탭 (Multi-tab):} 여러 쿼리를 동시에 열고 작업할 수 있습니다.
        \item \textbf{AI 어시스턴트 (AI Assistant):} 자연어로 질문하거나 코드 생성을 요청하여 쿼리 작성을 돕습니다.
        \item \textbf{SQL 웨어하우스 연결:} 쿼리 실행을 위해 SQL 엔드포인트(컴퓨팅 자원)에 연결합니다.
    \end{itemize}
    \item \textbf{SQL 웨어하우스 (SQL Warehouse):}
    \begin{itemize}
        \item 쿼리 실행을 위한 컴퓨팅 자원입니다.
        \item \textbf{서버리스 (Serverless):} Databricks 환경에서 VM(가상 머신)이 미리 준비되어 있어, 요청 시 즉시 컴퓨팅을 제공합니다. 비활성 상태가 10분 지속되면 자동으로 중지됩니다.
        \item \textbf{연결 세부 정보:} Tableau, PowerBI, dbt, Python, Java, NodeJS, Go 등 다양한 외부 BI 도구 및 애플리케이션이 JDBC/ODBC 연결을 통해 SQL 웨어하우스에 연결할 수 있습니다.
    \end{itemize}
\end{itemize}

\subsection{쿼리 프로파일 분석}
쿼리 프로파일은 SQL 쿼리의 성능을 이해하고 최적화하는 데 필수적인 정보를 제공합니다.

\begin{summarybox}[title=핵심 요약: 쿼리 성능 진단서]
쿼리 프로파일은 SQL 쿼리가 어떻게 실행되었는지, 얼마나 많은 시간과 메모리를 사용했는지, 어떤 단계에서 병목 현상이 발생했는지 등을 상세하게 보여주는 '진단서'입니다. 이를 통해 쿼리 성능을 개선하고 시스템 자원 사용을 최적화할 수 있습니다.
\end{summarybox}

\begin{enumerate}
    \item \textbf{쿼리 기록 (Query History) 접근:} Databricks SQL 인터페이스에서 'Query History' 섹션으로 이동합니다.
    \item \textbf{쿼리 필터링:} 지난 7일간의 쿼리, 특정 SQL 웨어하우스에서 실행된 쿼리, 1시간 이상 소요된 쿼리, 실패한 쿼리 등으로 필터링할 수 있습니다.
    \item \textbf{실패한 쿼리 분석:}
    \begin{itemize}
        \item 실패한 쿼리를 클릭하여 실패 원인(예: 권한 거부)을 확인합니다.
        \item 쿼리 실행 시간, 대기 시간 등을 파악하여 문제 해결에 활용합니다.
    \end{itemize}
    \item \textbf{쿼리 프로파일 상세 보기:}
    \begin{itemize}
        \item 성공한 쿼리를 선택하여 'Query Profile' 탭을 확인합니다.
        \item \textbf{실행 시간 (Duration):} 쿼리 실행에 소요된 총 시간.
        \item \textbf{메모리 사용량 (Memory Peak):} 쿼리 실행 중 사용된 최대 메모리.
        \item \textbf{반환된 행 수 (Number of Rows Returned):} 쿼리 결과로 반환된 행의 수.
        \item \textbf{소스 (Sources):} 쿼리가 데이터를 가져온 소스.
        \textbf{전체 스캔 (Full Scan):} 테이블 전체를 스캔했는지 여부.
        \item \textbf{시각화:} 쿼리 실행 계획을 시각적으로 보여주어 병목 현상이나 비효율적인 부분을 쉽게 파악할 수 있습니다.
    \end{itemize}
    \item \textbf{최적화 활용:} 쿼리 프로파일을 분석하여 인덱스 추가, 쿼리 재작성, 데이터 파티셔닝 등 쿼리 최적화 방안을 모색합니다.
\end{enumerate}

\subsection{경고(Alerts) 설정}
Databricks SQL에서 경고를 설정하여 데이터 파이프라인의 문제나 비즈니스 지표의 변화를 실시간으로 감지할 수 있습니다.

\begin{summarybox}[title=핵심 요약: 데이터 이상 감지 시스템]
경고는 특정 쿼리 결과가 미리 정의된 조건을 만족할 때 자동으로 알림을 보내는 '데이터 이상 감지 시스템'입니다. 이를 통해 데이터 품질 문제나 비즈니스 성과 저하를 조기에 파악하고 대응할 수 있습니다.
\end{summarybox}

\begin{enumerate}
    \item \textbf{경고 생성:} Databricks SQL 인터페이스에서 'Alerts' 섹션으로 이동하여 새로운 경고를 생성합니다.
    \item \textbf{쿼리 선택:} 경고를 모니터링할 SQL 쿼리를 선택합니다.
    \item \textbf{조건 정의:} 쿼리 결과에 대한 조건을 정의합니다 (예: '쿼리가 반환하는 행 수가 100개 미만일 때').
    \item \textbf{알림 빈도 설정:} 경고가 발생했을 때 알림을 보낼 빈도를 설정합니다 (예: 처음 발생 시, 매번 발생 시).
    \item \textbf{알림 대상 설정:} 알림을 받을 채널을 지정합니다 (예: Slack 채널, PagerDuty, 이메일).
    \item \textbf{활용 예시:}
    \begin{itemize}
        \item 데이터 파이프라인에서 예상보다 적은 데이터가 로드될 때 경고.
        \item 특정 비즈니스 지표(예: 매출)가 임계값 이하로 떨어질 때 경고.
        \item 중요한 데이터셋에 대한 비정상적인 접근 시도 감지.
    \end{itemize}
\end{enumerate}

\newpage
\section{체크리스트}

이 섹션은 강의 내용을 효과적으로 학습하고, 시험을 준비하며, Databricks 플랫폼을 활용하는 데 도움이 되는 점검표를 제공합니다.

\begin{summarybox}[title=학습 및 복습 체크리스트]
\begin{itemize}
    \item \textbf{데이터 저장소 이해:} 데이터 레이크, 데이터 웨어하우스, 레이크하우스의 정의, 장단점, 사용 사례를 명확히 구분할 수 있는가?
    \item \textbf{아키텍처 패턴:} 메달리온 아키텍처(Bronze, Silver, Gold)의 각 계층 역할과 중요성을 설명할 수 있는가?
    \item \textbf{데이터 관리 문제:} 데이터 사일로와 데이터 스웜프의 차이점과 발생 원인을 이해하는가?
    \item \textbf{현대 데이터 아키텍처:} 데이터 메시와 데이터 패브릭의 개념과 차이점을 설명할 수 있는가?
    \item \textbf{BI 기본:} 비즈니스 인텔리전스(BI)의 정의, 목표, 진화 과정을 이해하는가?
    \textbf{BI vs. BA:} BI와 BA의 주요 질문, 분석 유형, 목표 차이를 설명할 수 있는가?
    \item \textbf{BI 프로세스:} 데이터 수집부터 의사결정까지의 BI 프로세스 단계를 설명할 수 있는가?
    \item \textbf{BI 분석가 역할:} BI 분석가의 주요 기술, 사용 데이터 계층, 역할에 대해 이해하는가?
    \item \textbf{데이터 웨어하우스 용어:} 카탈로그, 스키마, 테이블, 키, 뷰, 구체화된 뷰 등 핵심 용어를 설명할 수 있는가?
    \item \textbf{데이터 모델링:} 3NF, 차원 모델링(스타/스노우플레이크), 데이터 볼트의 개념과 적용 시점을 이해하는가?
    \item \textbf{Databricks 플랫폼:} Databricks 플랫폼의 전반적인 아키텍처(소스, 수집, 변환, 쿼리, 스토리지, 거버넌스 등)를 이해하는가?
    \item \textbf{Databricks SQL 기능:} SQL 편집기, 대시보드, Genie, AI 쿼리/함수, 경고, 쿼리 프로파일 등 주요 기능을 설명할 수 있는가?
    \textbf{페더레이션 vs. ETL:} 데이터 페더레이션과 ETL의 차이점(특히 데이터 소유권)을 명확히 설명할 수 있는가?
    \item \textbf{BI 핵심 KPI:} 동시성(Concurrency)과 확장성(Scaling)이 BI 시스템에서 왜 중요한지 설명할 수 있는가?
    \item \textbf{실습 내용:} NYC Taxi 대시보드, Workspace Usage 대시보드, Genie 사용법, SQL 편집기, SQL 웨어하우스 연결, 쿼리 기록 및 프로파일 분석 방법을 이해하는가?
\end{itemize}
\end{summarybox}

\begin{summarybox}[title=시험 대비 전략 체크리스트]
\begin{itemize}
    \item \textbf{개념 정의 암기:} 모든 주요 용어의 정의와 원어를 정확히 암기했는가?
    \textbf{비교 분석:} 유사한 개념(예: Data Lake vs. DW vs. Lakehouse, BI vs. BA, Data Mesh vs. Fabric, Federation vs. ETL) 간의 차이점을 표로 정리하고 설명할 수 있는가?
    \item \textbf{단계별 프로세스:} BI 프로세스, 메달리온 아키텍처 등 단계별로 설명되는 내용을 순서대로 나열하고 각 단계의 역할을 설명할 수 있는가?
    \item \textbf{Databricks 기능 숙지:} Databricks SQL의 주요 기능들이 어떤 문제를 해결하고 어떤 이점을 제공하는지 설명할 수 있는가? (예: AI 쿼리/함수, Genie, 서버리스)
    \item \textbf{예시 활용:} 각 개념에 대한 현실적인 예시나 비유를 들어 설명할 수 있는가?
    \item \textbf{Q\&A 복습:} 강의 중 나왔던 질문과 답변들을 다시 확인하여 초심자가 혼동할 만한 포인트를 파악했는가?
\end{itemize}
\end{summarybox}

\newpage
\section{FAQ (자주 묻는 질문)}

초심자가 강의 내용을 학습하며 혼동할 수 있는 질문들을 Q\&A 형식으로 정리했습니다.

\begin{summarybox}[title=FAQ: 궁금증 해소]
\begin{itemize}
    \item \textbf{Q: 데이터 페더레이션(Data Federation)과 ETL의 가장 큰 차이점은 무엇인가요?}
    \item \textbf{A:} 가장 큰 차이점은 \textbf{데이터 소유권}입니다. ETL은 데이터를 원본에서 추출하여 변환한 후 데이터 레이크나 웨어하우스로 로드하여 \textbf{새로운 시스템이 데이터의 소유권을 갖게 됩니다}. 반면 페더레이션은 데이터를 물리적으로 이동시키지 않고 \textbf{원본 시스템에 소유권을 유지한 채} 읽기 전용으로 데이터를 통합하여 쿼리합니다. 페더레이션은 주로 참조 데이터나 소규모 데이터 통합에 사용되며, 대규모 데이터의 저지연/고처리량 분석에는 ETL이 더 적합합니다.

    \item \textbf{Q: Databricks SQL 대시보드를 사용하는 것이 PowerBI나 Tableau 같은 전문 BI 도구보다 어떤 장점이 있나요?}
    \item \textbf{A:} Databricks SQL 대시보드의 가장 큰 장점은 \textbf{라이선스 비용이 없다는 점}과 \textbf{레이크하우스 데이터에 대한 직접적인 접근성}입니다. PowerBI나 Tableau는 매우 정교한 시각화와 사용자 정의 기능을 제공하지만, 별도의 라이선스 비용이 발생합니다. Databricks SQL 대시보드는 운영 대시보드나 빠른 프로토타이핑에 적합하며, 레이크의 최신 데이터에 직접 연결하여 별도의 데이터 웨어하우스를 거치지 않고도 고성능 분석을 할 수 있습니다. 외부 사용자에게 임베디드 자격 증명으로 쉽게 게시할 수 있는 것도 큰 장점입니다.

    \item \textbf{Q: Genie는 Chat GPT와 어떻게 다른가요? Genie의 답변을 신뢰할 수 있나요?}
    \item \textbf{A:} Genie는 Chat GPT와 달리 \textbf{인터넷의 일반적인 지식이 아닌, 사용자가 제공한 대시보드의 메타데이터와 데이터셋만을 기반으로 답변}합니다. 따라서 '환각(Hallucination)' 현상이 거의 없으며, 질문이 불분명할 경우 답변을 거부하거나 명확화를 요청합니다. Genie는 데이터셋 내의 정보만을 활용하므로, 답변의 신뢰성이 높습니다. 또한, Genie가 생성한 SQL 쿼리를 직접 확인하여 답변의 근거를 검증할 수 있습니다.

    \item \textbf{Q: Databricks에서 대시보드를 게시할 때 발생하는 비용은 누가 부담하나요?}
    \item \textbf{A:} 대시보드를 '임베디드 자격 증명(Embedded Credentials)'으로 게시하여 외부 사용자가 접근할 경우, 대시보드가 새로 고쳐지거나 캐시된 결과가 만료되어 새로운 쿼리가 실행될 때 발생하는 \textbf{컴퓨팅 비용은 대시보드를 게시한 사용자(계정)가 부담}합니다. 서버리스 환경에서는 기본 컴퓨팅 인프라 비용과 Databricks 플랫폼 사용 비용이 통합되어 청구됩니다.

    \item \textbf{Q: 데이터 계보(Data Lineage)와 감사(Auditing)는 어떻게 다른가요?}
    \item \textbf{A:} \textbf{데이터 계보(Lineage)}는 데이터가 원본에서 최종 목적지까지 어떻게 흘러가고 어떤 변환을 거쳤는지 추적하는 것입니다. 데이터의 출처와 변환 과정을 이해하는 데 중요합니다. 반면 \textbf{감사(Auditing)}는 누가 어떤 테이블에 언제 접근했는지, 접근이 성공했는지 실패했는지 등을 기록하는 것입니다. 이는 민감한 데이터에 대한 접근을 모니터링하고, 규제 준수 및 내부 보안 위협 감지에 사용됩니다.
\end{itemize}
\end{summarybox}

\newpage
\section{빠르게 훑어보기 (1페이지 요약)}

\begin{summarybox}[title=CSCI103 Lecture 6 핵심 요약]
\textbf{데이터 저장소의 진화:} 스프레드시트 $\rightarrow$ DW $\rightarrow$ MPP $\rightarrow$ NoSQL $\rightarrow$ Hadoop $\rightarrow$ Data Lake $\rightarrow$ Data Mesh/Fabric $\rightarrow$ Lakehouse. AI/LLM으로 가속화 중.

\textbf{데이터 레이크 vs. DW vs. 레이크하우스:}
\begin{itemize}
    \item \textbf{Data Lake:} 모든 데이터 유형, 스키마 온 리드, 저렴한 스토리지, 유연성. (단점: 거버넌스 부족 시 스웜프)
    \item \textbf{Data Warehouse:} 정형 데이터, 스키마 온 라이트, 고성능 쿼리, 낮은 지연 시간. (단점: 데이터 양 제한, 비정형 처리 어려움)
    \item \textbf{Lakehouse:} Data Lake의 유연성 + DW의 성능/거버넌스 결합. 모든 데이터 유형, 고성능 SQL, 비용 효율적.
\end{itemize}

\textbf{메달리온 아키텍처:} 데이터 품질에 따른 3단계 계층화.
\begin{itemize}
    \item \textbf{Bronze:} 원본 데이터 (진실의 원천)
    \item \textbf{Silver:} 정제 및 통합 데이터 (데이터 과학자용)
    \item \textbf{Gold:} 분석 준비 완료 데이터 (비즈니스 사용자용)
\end{itemize}

\textbf{BI의 이해:} 데이터를 통해 비즈니스 의사결정 개선.
\begin{itemize}
    \item \textbf{BI (기술적 분석):} "무엇이 일어났는가?" (과거 지향)
    \item \textbf{BA (예측/처방적 분석):} "왜 일어났고, 다음에 무엇을 해야 하는가?" (미래 지향)
    \item \textbf{BI 프로세스:} 수집 $\rightarrow$ 조직 $\rightarrow$ 분석 $\rightarrow$ 공유 $\rightarrow$ 활용.
\end{itemize}

\textbf{Databricks 플랫폼 핵심:}
\begin{itemize}
    \item \textbf{Unity Catalog:} 통합 데이터 거버넌스 (메타데이터, 계보, 접근 제어).
    \item \textbf{Databricks SQL (DBSQL):} BI/DW를 위한 SQL 인터페이스. Photon 엔진으로 고성능.
    \item \textbf{페더레이션:} 데이터 소유권 없이 외부 데이터 통합 쿼리 (ETL과 구분).
    \item \textbf{Genie:} 대시보드에서 자연어로 질문하고 답변 얻는 AI 기능.
    \item \textbf{AI 쿼리/함수:} SQL 내에서 LLM 호출 (감성 분석, 분류, 정보 추출 등).
    \textbf{서버리스:} 즉시 컴퓨팅 제공, 높은 동시성 및 확장성.
    \item \textbf{대시보드:} 쉽게 생성, 게시, 경고 설정 가능.
\end{itemize}

\textbf{BI 분석가:} SQL이 핵심 기술, Gold/Silver 계층 데이터 사용, 셀프 서비스 BI 지원.

\textbf{핵심 KPI:}
\begin{itemize}
    \item \textbf{동시성 (Concurrency):} 동시에 처리 가능한 쿼리 수.
    \item \textbf{확장성 (Scaling):} 데이터/사용자 증가에 따른 성능 유지 능력.
\end{itemize}
\end{summarybox}

\end{document}
